\FloatBarrier
\section{Double-Dot Charge Stability Diagrams}
\label{app:charge-stability}
\label{sec:charge-stability-measurements}

Section~\ref{sec:double-dot-loading} uses the charge regions in
Fig.~\ref{fig:measurement-honeycomb} to choose voltages for preparing
one electron on each dot. This appendix derives their boundaries
from the electron-number condition and relates them to measured
charge-transition lines.

\subsection{Charge regions and their boundaries}
\label{sec:measurement-double-dot-csd}

For two dots, the desired result is a map of the charge configuration
$(N_1,N_2)$ as a function of the two plunger voltages $V_1,V_2$.
This locates settings for loading each dot and for transferring an
electron between dots. Keep both reservoirs at a common
$\mu_{\mathrm R}$ and first consider weak interdot tunneling, so
that the occupations label approximately distinct configurations.
Let $E_{N_1,N_2}(\mathbf V)$ be the lowest energy of each
configuration before tunneling mixes it with the others. The
equilibrium pair minimizes
\begin{equation}
 \Omega_{N_1,N_2}(\mathbf V)
       =E_{N_1,N_2}(\mathbf V)-\mu_{\mathrm R}(N_1+N_2).
 \label{eq:measurement-double-dot-grand-energy}
\end{equation}
This extends the electron-number construction in Section~\ref{sec:electron-number-control}
to two separately resolved dots.

A boundary occurs when two minimizing configurations have equal
$\Omega$. Adding an electron from a reservoir to dot 1 gives
\begin{equation}
 \mu_1(N_1+1,N_2)\equiv
 E_{N_1+1,N_2}-E_{N_1,N_2}=\mu_{\mathrm R}.
 \label{eq:measurement-dot1-boundary}
\end{equation}
There is an analogous condition for dot 2. An interdot transfer
conserves total electron number, so the reservoir term cancels:
\begin{equation}
 E_{N_1+1,N_2-1}=E_{N_1,N_2}.
 \label{eq:measurement-interdot-boundary}
\end{equation}
These are different physical events: the first changes the total
charge, while the second redistributes it.

Repulsion between dots makes the addition threshold of one depend
on the occupation of the other. A line for adding to dot 1 therefore
shifts when dot 2 gains an electron. Together with interdot-transfer
segments, the shifted lines commonly form hexagonal cells, called
the \emph{honeycomb diagram}~\cite{VanDerWiel2003}.
Figure~\ref{fig:measurement-honeycomb} shows one cell and its neighbors.
Its polygon is a region in voltage space, not the spatial shape of
a dot. Varying orbital spacings, interactions, or gate responses
changes the cell sizes and slopes.

The slopes follow directly from the boundary equations. Locally
write the change in the relevant addition threshold as
\begin{equation}
 \delta\mu_i=-e(\alpha_{i1}^{(\mu)}\delta V_1
                  +\alpha_{i2}^{(\mu)}\delta V_2),
 \qquad
 \left.\frac{dV_2}{dV_1}\right|_{\mu_i=\mu_{\mathrm R}}
      =-\frac{\alpha_{i1}^{(\mu)}}{\alpha_{i2}^{(\mu)}}.
 \label{eq:measurement-csd-slope}
\end{equation}
Here $\alpha_{ig}^{(\mu)}=-(1/e)\partial\mu_i/\partial V_g$
is a dimensionless addition-threshold lever arm. The superscript
distinguishes it from a single-orbital lever arm in
Eq.~\eqref{eq:device-lever-arm}. A slope gives their ratio, but
not an absolute energy scale.

Charge sensing can reveal the boundary lines throughout the map.
Near-zero-bias sequential transport through the double dot is
concentrated near the triple points, where three charge regions
meet; Section~\ref{sec:measurement-bias-triangles} explains this
transport condition.
To obtain the map by charge sensing, fix the reservoir and barrier-gate voltages.
Sweep $V_1$ at a sequence of fixed $V_2$ values and record the
charge-sensor signal at each pair. The resulting two-dimensional
record is a \emph{charge stability diagram} (CSD). Its changes can
be compared with the boundaries just derived. Establish an empty
reference region, then count loading transitions to assign absolute
electron numbers. If tunneling becomes too slow, an occupied region
can appear empty on the measurement timescale; the waiting-time
test below is therefore part of this assignment.

\subsection{Extracting transition lines from a sensor record}
\label{sec:measurement-charge-scan}

We want to make the boundaries between charge regions easy to
identify in the measured voltage map. Recall the local sensor model:
\begin{equation*}
 s(t)=s_{\mathrm{bg}}(\mathbf V)
         +\sum_i\kappa_i N_i(t)+\xi(t).
 \tag{\ref*{eq:measurement-sensor-signal}}
\end{equation*}
Within each region, the fixed occupations give a fixed charge
contribution $\sum_i\kappa_iN_i$, while $s_{\mathrm{bg}}$ varies
smoothly with gate voltage. When the scan crosses a boundary where
dot $i$ gains an electron, the extra charge shifts the sensor's
addition energy and hence its current. The change $N_i\to N_i+1$
changes the charge contribution by $\kappa_i$, producing a step on
top of the smooth background. At finite temperature, the mean
signal changes smoothly between the two levels.

Differentiating with respect to the scanned gate voltage turns
this step into a peak or dip, highlighting the boundary as a line
in the two-dimensional map; within each region, the fixed charge
contributes no derivative. For zero-mean detector noise, averaging
the signal and differentiating with respect to $V_g$, with the
other voltages fixed, gives
\begin{equation}
 \frac{\partial\overline s}{\partial V_g}
   =\frac{\partial s_{\mathrm{bg}}}{\partial V_g}
      +\sum_i\kappa_i
          \frac{\partial\langle\widehat N_i\rangle}{\partial V_g}.
 \label{eq:measurement-csd-response}
\end{equation}

Clear steps can also be located directly in the original signal;
differentiation is optional, and numerical differentiation can
amplify measurement noise.

An interdot transfer from dot 2 to dot 1 changes the signal by
approximately $\kappa_1-\kappa_2$. Its line can therefore be faint
when the two sensor couplings are similar.

The voltage derivative can also be measured by applying a small
gate modulation. Let $S(V_g)$ be the mean sensor signal at a fixed
gate voltage $V_g$, with all other voltages fixed. Apply
$V_g(t)=V_g^0+v_{\mathrm{ac}}\cos\omega_m t$ slowly enough that the
charge follows the modulation. Omitting detector noise for the
derivation, $s(t)=S[V_g(t)]$. A Taylor expansion about $V_g^0$ gives
\[
 s(t)=S(V_g^0)
       +v_{\mathrm{ac}}S'(V_g^0)\cos\omega_m t
       +O(v_{\mathrm{ac}}^2).
\]
Multiply both sides by $\cos\omega_m t$ and integrate over a time
$\tau_{\mathrm{av}}$ containing an integer number of modulation
periods. Keeping the first-order terms gives
\[
 \begin{aligned}
 \frac{1}{\tau_{\mathrm{av}}}
   \int_0^{\tau_{\mathrm{av}}}s(t)\cos\omega_m t\,dt
 &\simeq
 \frac{S(V_g^0)}{\tau_{\mathrm{av}}}
   \int_0^{\tau_{\mathrm{av}}}\cos\omega_m t\,dt
 \\
 &\quad+
 \frac{v_{\mathrm{ac}}S'(V_g^0)}{\tau_{\mathrm{av}}}
   \int_0^{\tau_{\mathrm{av}}}\cos^2\omega_m t\,dt
 \\
 &=\frac{v_{\mathrm{ac}}}{2}S'(V_g^0).
 \end{aligned}
\]
The two integrals on the right are $0$ and
$\tau_{\mathrm{av}}/2$, respectively. Writing the time average as
$\overline f=\tau_{\mathrm{av}}^{-1}
\int_0^{\tau_{\mathrm{av}}}f(t)\,dt$, we obtain
\begin{equation}
 \frac{2}{v_{\mathrm{ac}}}
    \overline{s(t)\cos\omega_m t}
       \simeq \left.\frac{dS}{dV_g}\right|_{V_g^0}.
 \label{eq:measurement-lock-in}
\end{equation}
An instrument that performs this reference multiplication and
averaging is a \emph{lock-in amplifier}. Detector noise adds
uncertainty to this estimate. The Taylor approximation requires a
modulation amplitude small compared with the voltage scale over
which the response varies. If the charge cannot follow the
modulation, the measured amplitude and phase also contain dynamics.

To locate the charge-transition boundaries experimentally, we
construct a map of $\partial\overline s/\partial V_1$ or
$\partial\overline s/\partial V_2$ over the $(V_1,V_2)$ plane.
There are two ways to obtain this map.

For numerical differentiation, scan $V_1$ at a sequence of fixed
$V_2$ values and record the mean sensor signal at each voltage
pair. Divide the signal difference between neighboring samples
by their voltage separation to estimate the corresponding
derivative. Plotting these estimates over $(V_1,V_2)$ highlights
the charge-transition boundaries.

For lock-in measurement, scan the same voltage pairs, but apply
a small modulation to $V_1$ or $V_2$ at each pair. Record the
lock-in response and use Eq.~\eqref{eq:measurement-lock-in} to
obtain the corresponding derivative directly. Choose a modulation
small compared with the feature width and slow enough for charge
to follow it.

To map equilibrium charge states, allow the electrons to equilibrate
at each gate-voltage setting before recording. Otherwise a charge
transition can be missed because an electron has not yet entered.
For an initially empty accessible state with constant loading rate
$\Gamma_{\mathrm{in}}$ and negligible unloading, the probability
of loading during a waiting time $\tau$ is
\begin{equation}
 P_{\mathrm{load}}(\tau)=1-e^{-\Gamma_{\mathrm{in}}\tau}.
 \label{eq:measurement-loading-probability}
\end{equation}
The wait must therefore be long compared with
$\Gamma_{\mathrm{in}}^{-1}$, as well as any relevant relaxation
times. To check whether insufficient waiting caused a missing line,
repeat that region with a longer wait at each voltage setting.
A line restored by allowing more time to load
reveals a failure to equilibrate in the faster scan, an important
test in the last-electron regime~\cite{Thalakulam2010}.

\subsection{Detecting transport cycles: triangles and triple points}
\label{sec:measurement-bias-triangles}

We want to identify the gate-voltage pairs that support a current
through two dots in series. First consider positive finite bias,
$\mu_{\mathrm S}-\mu_{\mathrm D}=eV_b>0$. A forward cycle requires
an electron to enter dot 1 from the source, transfer to dot 2,
and leave into the drain. Denote the configurations by
$A=(N_1,N_2)$, $B=(N_1+1,N_2)$, and $C=(N_1,N_2+1)$.
The cycle is $A\to B\to C\to A$. Define the addition energies
$\widetilde\mu_1=E_B-E_A$ and
$\widetilde\mu_2=E_C-E_A$, with all energies evaluated at the
chosen bias. At zero temperature and negligible tunnel broadening,
entry, downhill transfer, and exit require
\begin{equation}
 \mu_{\mathrm S}\geq\widetilde\mu_1
       \geq\widetilde\mu_2\geq\mu_{\mathrm D}.
 \label{eq:measurement-triangle-window}
\end{equation}
The first inequality permits loading, the middle one permits
interdot transfer from $B$ to $C$ without absorbing energy, since
$E_C-E_B=\widetilde\mu_2-\widetilde\mu_1$, and the last permits
unloading.
In the $(\widetilde\mu_1,\widetilde\mu_2)$ plane they enclose a
triangle. Approximately linear gate responses map it to a triangle
in $(V_1,V_2)$, whose extent grows with $eV_b$.

The interior requires a way to release the interdot energy
difference, for example by emitting a lattice vibration, or phonon.
Purely elastic transfer between two sharp levels is resonant along
$\widetilde\mu_1=\widetilde\mu_2$ instead. Excited states create
additional resonances, and spin selection can block an otherwise
allowed cycle. A bias triangle is therefore an allowed transport
region, not necessarily a uniformly bright patch~\cite{VanDerWiel2003}.

Now reduce the bias toward zero. As both reservoir chemical
potentials approach a common $\mu_{\mathrm R}$,
Eq.~\eqref{eq:measurement-triangle-window} forces
\[
 \widetilde\mu_1=\widetilde\mu_2=\mu_{\mathrm R}.
\]
Because $B$ and $C$ each contain one more electron than $A$,
their grand-energy differences are
$\Omega_B-\Omega_A=\widetilde\mu_1-\mu_{\mathrm R}$ and
$\Omega_C-\Omega_A=\widetilde\mu_2-\mu_{\mathrm R}$. Hence
\begin{equation}
 \Omega_A=\Omega_B=\Omega_C.
 \label{eq:measurement-triple-point}
\end{equation}
These two independent equalities select isolated points in the
two-dimensional voltage plane: the finite-bias triangles shrink
to \emph{triple points}. The nearby complementary cycle through
$(N_1+1,N_2+1)$ gives the other member of a triple-point pair.
Near zero bias, the current is confined to small neighborhoods
of these points, with finite temperature broadening the features.
At exactly zero DC bias, the net current vanishes; it is the
linear conductance that has peaks near the triple points.

To locate these transport regions experimentally, connect the two
dots in series between source and drain, choose a positive finite
bias, and record the current while scanning the two plunger
voltages. Keep the reservoir voltages fixed during each map.
Repeat at successively smaller biases to follow the triangles
as they shrink toward the triple points. Reversing the bias probes
transport in the opposite direction. Compare these maps with
charge-sensor boundaries acquired at the corresponding settings.
A sensor line can identify a charge transition even where no
complete transport cycle is resolved. The triangle edges, rather
than their brightness alone, are the features to compare with
Eq.~\eqref{eq:measurement-triangle-window}.
